\documentclass[a4paper,11pt,openany,extrafontsizes]{memoir}

\usepackage{fontspec}

\setmainfont[Numbers=OldStyle]{Linux Libertine O}
\setsansfont[Numbers=OldStyle]{Linux Biolinum O}
\setmonofont[Scale=0.83]{Inconsolata}

\usepackage{polyglossia}
\setdefaultlanguage[variant=british]{english}

\usepackage{lipsum}

\usepackage{graphicx}
\usepackage[dvipsnames]{xcolor}
\usepackage{wrapfig}
\usepackage{subcaption}
\usepackage{lettrine}

\usepackage{amssymb, amsmath}
\usepackage{amsthm}

\theoremstyle{plain}
\newtheorem{thm}{Theorem}[chapter]
\newtheorem{lem}[thm]{Lemma}
\newtheorem{cor}[thm]{Corollary}
\newtheorem{prop}[thm]{Proposition}
\theoremstyle{definition}
\newtheorem{defn}{Definition}[chapter]
\newtheorem{expl}{Example}[chapter]
\theoremstyle{remark}
\newtheorem*{rem}{Remark}
\newtheorem*{note}{Note}
\newtheorem*{notation}{Notation}

\usepackage{pdfpages}

\usepackage{microtype}

%% Propriétés du document PDF
\usepackage[unicode,colorlinks=true]{hyperref}

\hypersetup{
  pdfauthor={Dimitri Lozeve},
  pdftitle={Topological Data Analysis of Temporal Networks},
  pdfsubject={A dissertation submitted un partial fulfilment of the degree of Master of Science in Applied Statistics},
  pdfkeywords={dissertation,stats,msc,tda,networks},
  pdfpagemode=UseOutlines,
  pdfpagelayout=TwoColumnRight,
  linkcolor=MidnightBlue,
  filecolor=MidnightBlue,
  urlcolor=MidnightBlue,
  citecolor=Green
}

%% Pour la classe memoir /!\

%% Marges
\setlrmarginsandblock{3cm}{3cm}{*}
\setulmarginsandblock{3cm}{3cm}{*}
\checkandfixthelayout%

%% Numérotation des divisions logiques
\setsecnumdepth{subsection}
\maxsecnumdepth{subsection}

%% Profondeur de la ToC
\settocdepth{subsection}
\maxtocdepth{subsection}

%% Style des titres des divisions logiques
\setsecheadstyle{\Large\scshape}
\setsubsecheadstyle{\large\scshape}

%% Abstract
\abstractintoc%
\renewcommand{\abstractnamefont}{\normalfont\large\scshape}
\renewcommand{\abstracttextfont}{\normalfont\normalsize}

%% épigraphes
\setlength{\epigraphwidth}{0.5\textwidth}
\epigraphtextposition{flushleftright}

%% Couleurs
%\definecolor{purpletouch}{RGB}{103,30,117}
\definecolor{bleux}{RGB}{0,62,92}

\author{Dimitri Lozeve}
\date{September 2018}
\title{Topological Data Analysis of Temporal Networks}







%%% Local Variables:
%%% mode: latex
%%% TeX-master: "dissertation"
%%% End:


\usepackage[firstpage]{draftwatermark}


\begin{document}

\pagestyle{plain}
\tightlists%

\begin{titlingpage}
  \begin{center}
    \vspace{1cm}
    \textsf{\Huge{University of Oxford}}\\
    \vspace{1cm}
    \includegraphics[scale=.8]{Stats_Logo.png}\\
    \vspace{2cm}
    \Huge{\thetitle}\\
    \vspace{2cm}
    \large{by\\[14pt]\theauthor\\[8pt]St Catherine's College}\\
    % \vspace{2.2cm}
    \vfill
    \large{A dissertation submitted in partial fulfilment of the degree of Master of Science in Applied Statistics}\\
    \vspace{.5cm}
    \large{\emph{Department of Statistics, 24--29 St Giles,\\Oxford, OX1 3LB}}\\
    \vspace{1cm}
    \large{\thedate}
  \end{center}
\end{titlingpage}

%\chapterstyle{hangnum}
%\chapterstyle{ell}
%\chapterstyle{southall}
\chapterstyle{wilsondob}

\frontmatter

\cleardoublepage%

\chapter*{Declaration of authorship}

\emph{This my own work (except where otherwise indicated).}\\[2cm]

\begin{center}
  Date \hspace{.5\linewidth} Signature
\end{center}


\cleardoublepage%

\begin{abstract}
  Abstract here
\end{abstract}

\cleardoublepage%

\chapter*{Acknowledgements}%
\label{cha:acknowledgements}

Thank you!

\cleardoublepage%

\tableofcontents*
\listoffigures*
\listoftables*

\clearpage

\mainmatter%

\chapter{Introduction}%
\label{cha:introduction}


\chapter{Graphs and Temporal Networks}%
\label{cha:temporal-networks}

\section{Definition and basic properties}%
\label{sec:defin-basic-prop}

In this section, we will introduce the notion of temporal networks or
graphs. This is a complex notion, with many concurrent definitions and
interpretations. First, we restate the standard definition of a
non-temporal, static graph.

\begin{defn}[Graph]
  A \emph{graph} is a couple $G = (V, E)$, where $V$ is a finite set
  of \emph{nodes} (or \emph{vertices}), and $E \subseteq V\times V$ is
  a set of \emph{edges}. A \emph{weighted graph} is defined by
  $G = (V, E, w)$, where $w : E\mapsto \mathbb{R}_+$ is fcalled the
  \emph{weight function}.
\end{defn}

We also define some basic concepts that will be needed later on to
build simplicial complexes on graphs.

\begin{defn}[Clique]
  A \emph{clique} is a set of nodes where each pair is connected. That
  is, a clique $C$ of a graph $G = (V,E)$ is a subset of $V$ such that
  $\forall i,j\in C, i \neq j \implies (i,j)\in E$. A clique is said
  to be \emph{maximal} if it cannot be augmented by any node.
\end{defn}

Temporal networks are defined in the more general framework of
\emph{multilayer networks}~\cite{kivela_multilayer_2014}. However,
this definition is much too general for our simple applications, and
we restrict ourselves to edge-centric time-varying
graphs~\cite{casteigts_time-varying_2012}. In this model, the set of
nodes is fixed and doesn't change over time, whereas edges can appear
or disappear at different timestamps.

\begin{defn}[Temporal network]
  A \emph{temporal network} (or graph) is a tuple
  $G = (V, E, \mathcal{T}, \rho)$, where:
  \begin{itemize}
  \item $V$ is a finite set of nodes,
  \item $E\subseteq V\times V$ is a set of edges,
  \item $\mathbb{T}$ is the \emph{temporal domain} (often taken as
    $\mathbb{N}$ or $\mathbb{R}_+$), and
    $\mathcal{T}\subseteq\mathbb{T}$ is the \emph{lifetime} of the
    network,
  \item $\rho: E\times\mathcal{T}\mapsto\{0,1\}$ is the \emph{presence
      function}, which determines whether an edge is present in the
    network at each timestamp.
  \end{itemize}
  The \emph{available dates} of an edge are the set
  $\mathcal{I}(e) = \{t\in\mathcal{T}: \rho(e,t)=1\}$.
\end{defn}

Temporal networks can also have weighted edges. In this case, it is
possible to have constant weights (edges can only appear or disappear
over time, and always have the same weight), or time-varying
weights. In the latter case, we can set the domain of the presence
function to be $\mathbb{R}_+$ instead of $\{0,1\}$, where by
convention a zero weight corresponds to an absent edge.

\begin{defn}[Additive temporal network]
  A temporal network is said to be \emph{additive} if for all $e\in E$
  and $t\in\mathcal{T}$, if $\rho(e,t)=1$, then
  $\forall t'>t, \rho(e, t') = 1$. Edges can only be added to the
  network, never removed.
\end{defn}

\section{Examples of applications}%
\label{sec:exampl-appl}

\section{Network partitioning}%
\label{sec:network-partitioning}

Temporal networks are a very active research subject, leading to
multiple interesting problems. The additional time dimension adds a
significant layer of complexity that cannot be adequately treated by
the common methods on static graphs.

Moreover, data collection can lead to large amount of noise in
datasets. Combined with large dataset sized due to the huge number of
data points for each node in the network, temporal graphs cannot be
studied effectively in their raw form. Recent advances have been made
to fit network models to rich but noisy
data~\cite{newman_network_2018}, generally using some variation on the
expectation-maximization (EM) algorithm.

One solution that has been proposed to study such temporal data has
been to \emph{partition} the time scale of the network into a sequence
of smaller, static graphs, representing all the interactions during a
short interval of time. The approach consists in subdividing the
lifetime of the network in \emph{sliding windows} of a given length.
We can then ``flatten'' the temporal network on each time interval,
keeping all the edges that appear at least once (or adding their
weights in the case of weighted networks).

This partitioning is sensitive to two parameters: the length of each
time interval, and their overlap. Of those, the former is the most
important: it will define the \emph{resolution} of the study. If it is
too small, too much noise will be taken into account; if it is too
large, we will lose important information. There is a need to find a
compromise, which will depend on the application and on the task
performed on the network. In the case of a classification task to
determine periodicity, it will be useful to adapt the resolution to
the expected period: if we expect week-long periodicity, a resolution
of one day seems reasonable.

Once the network is partitioned, we can apply any statistical learning
task on the sequence of static graphs. In this study, we will focus on
classification of time steps. This can be used to detect periodicity,
outliers, or even maximise temporal communities.

%% TODO Talk about partitioning methods?

\chapter{Topological Data Analysis and Persistent Homology}%
\label{cha:tda-ph}

\section{Basic constructions}%
\label{sec:basic-constructions}

\subsection{Homology}%
\label{sec:homology}

Our goal is to understand the topological structure of a metric
space. For this, we can use \emph{homology}, which consists in
associating for a metric space $X$ and a dimension $i$ a vector space
$H_i(X)$. The dimension of $H_i(X)$ will give us the number of
$i$-dimensional components in $X$: the dimension of $H_0(X)$ is the
number of path-connected components in $X$, the dimension of $H_1(X)$
is the number of holes in $X$, and the dimension of $H_2(X)$ is the
number of voids.

Crucially, these vector spaces are robust to continuous deformation of
the underlying metric space (they are \emph{homotopy
  invariant}). However, computing the homology of an arbitrary metric
space can be extremely difficult. It is necessary to approximate it in
a structure that would be both combinatorial and topological in
nature.

\subsection{Simplicial Complexes}%
\label{sec:simplicial-complexes}

In order to understand the topological structure of a metric space, we
need a way to decompose it in smaller pieces which, when assembled,
conserve the overall organisation of the space. For this, we use a
structure called a \emph{simplicial complex}, which is a kind of
higher-dimensional generalization of graphs.

The building blocks of this representation will be \emph{simplices},
which are simply the convex hull of an arbitrary set of
points. Examples of simplices include single points, segments,
triangles, and tetrahedrons (in dimensions 0, 1,, 2, and 3
respectively).

\begin{defn}[Simplex]
  The \emph{$k$-dimensional simplex} $\sigma = [x_0,\ldots,x_k]$ is
  the convex hull of the set $\{x_0,\ldots,x_k\} \in \mathbb{R}^d$,
  where $x_0,\ldots,x_k$ are affinely independent. $x_0,\ldots,x_k$
  are called the \emph{vertices} of $\sigma$, and the simplices
  defined by the subsets of $\{x_0,\ldots,x_k\}$ are called the
  \emph{faces} of $\sigma$.
\end{defn}

\begin{figure}[ht]
  \centering
  \begin{subfigure}[b]{.3\linewidth}
    \centering
    \begin{tikzpicture}
      \tikzstyle{point}=[circle,thick,draw=black,fill=blue!30,%
      inner sep=0pt,minimum size=15pt]
      \node (a)[point] at (0,0) {a};
    \end{tikzpicture}
    \caption{Single vertex}
  \end{subfigure}%
  % 
  \begin{subfigure}[b]{.3\linewidth}
    \centering
    \begin{tikzpicture}
      \tikzstyle{point}=[circle,thick,draw=black,fill=blue!30,%
      inner sep=0pt,minimum size=15pt]
      \node (a)[point] at (0,0) {a};
      \node (b)[point] at (1.4,2) {b};
      
      \begin{scope}[on background layer]
        \draw[fill=blue!15] (a.center) -- (b.center) -- cycle;
      \end{scope}
    \end{tikzpicture}
    \caption{Segment}
  \end{subfigure}%
  % 
  \begin{subfigure}[b]{.3\linewidth}
    \centering
    \begin{tikzpicture}
      \tikzstyle{point}=[circle,thick,draw=black,fill=blue!30,%
      inner sep=0pt,minimum size=15pt]
      \node (a)[point] at (0,0) {a};
      \node (b)[point] at (1.4,2) {b};
      \node (c)[point] at (2.8,0) {c};
      
      \begin{scope}[on background layer]
        \draw[fill=blue!15] (a.center) -- (b.center) -- (c.center) -- cycle;
      \end{scope}
    \end{tikzpicture}
    \caption{Triangle}
  \end{subfigure}%
  % 
  \caption{Examples of simplices}%
  \label{fig:simplex}
\end{figure}


We then need a way to combine these basic building blocks meaningfully
so that the resulting object can adequately reflect the topological
structure of the metric space.

\begin{defn}[Simplicial complex]
  A \emph{simplicial complex} is a collection $K$ of simplices such
  that:
  \begin{itemize}
  \item any face of a simplex of $K$ is a simplex of $K$
  \item the intersection of two simplices of $K$ is either the empty
    set or a common face or both.
  \end{itemize}
\end{defn}

\begin{figure}[ht]
  \centering
  \begin{tikzpicture}
    \tikzstyle{point}=[circle,thick,draw=black,fill=blue!30,%
      inner sep=0pt,minimum size=10pt]
      \node (a)[point] {};
      \node (b)[point,above right=1.4cm and 1cm of a] {};
      \node (c)[point,right=2cm of a] {};
      \node (d)[point,above right=.4cm and 2cm of b] {};
      \node (e)[point,above right=.4cm and 2cm of c] {};
      \node (f)[point,below right=.7cm and 1.3cm of c] {};
      \node (g)[point,right=2cm of d] {};
      \node (h)[point,below right=.4cm and 1.5cm of e] {};
      
      \begin{scope}[on background layer]
        \draw[fill=blue!15] (a.center) -- (b.center) -- (c.center) -- cycle;
        \draw (b) -- (d) -- (g);
        \draw (c.center) -- (e.center) -- (f.center) -- cycle;
        \draw (d) -- (e) -- (h);
      \end{scope}

      \node (1)[point,right=2cm of g] {};
      \node (2)[point,above right=.5cm and 1cm of 1] {};
      \node (3)[point,below right=.5cm and 1cm of 2] {};
      \node (4)[point,below left=1cm and .3cm of 3] {};
      \node (5)[point,below right=1cm and .3cm of 1] {};
      \node (6)[point,below left=1cm and .1cm of 5] {};
      \node (7)[point,below right=1cm and .1cm of 4] {};
      \node (8)[point,below right=.7cm and .7cm of 6] {};

      \begin{scope}[on background layer]
        \draw[fill=green!15] (1.center) -- (2.center) -- (3.center) -- (4.center) -- (5.center) -- cycle;
        \draw (1) -- (4) -- (2) -- (5) -- (3) -- (1);
        \draw[fill=blue!15] (6.center) -- (7.center) -- (8.center) -- cycle;
        \draw (5) -- (6) -- (4) -- (7);
      \end{scope}
  \end{tikzpicture}
  \caption{Example of a simplicial complex, with two connected
    components, two 3-simplices, and one 5-simplex.}%
  \label{fig:simplical-complex}
\end{figure}

The notion of simplicial complex is closely related to that of a
hypergraph. The important distinction lies in the fact that a subset
of a hyperedge is not necessarily a hyperedge itself.

Using these definitions, we can define homology on simplicial
complexes. %% TODO add reference for more details/do it myself?

\subsection{Filtrations}%
\label{sec:filtrations}

If we consider that a simplicial complex is a kind of
``discretization'' of a metric space, we realise that there must be an
issue of \emph{scale}. For our analysis to be invariant under small
perturbations in the data, we need a way to find the optimal scale
parameter to capture the adequate topological structure, without
taking into account some small perturbations, nor ignoring some
important smaller features.

%% TODO rewrite using the Cech filtration as an example?

The ideal solution to these problems is to consider all scales at
once: this is the objective of \emph{filtered simplical complexes}.

\begin{defn}[Filtration]
  A \emph{filtered simplicial complex}, or simply a \emph{filtration},
  $K$ is a sequence ${(K_i)}_{i\in I}$ of simplicial complexes such
  that:
  \begin{itemize}
  \item for any $i, j \in I$, if $i < j$ then $K_i \subseteq K_j$,
  \item $\bigcup_{i\in I} K_i = K$.
  \end{itemize}
\end{defn}

\section{Persistent Homology}%
\label{sec:persistent-homology}

We can now compute the homology for each step in a filtration. This
leads to the notion of \emph{persistent
  homology}~\cite{carlsson_topology_2009,zomorodian_computing_2005},
which gives us all the information necessary to establish the
topological structure of the metric space at multiple scales.

\begin{defn}[Persistent homology]
  The \emph{$p$-th persistent homology} of a simplicial complex
  $K = {(K_i)}_{i\in I}$ is the pair
  $(\{H_p(K_i)\}_{i\in I}, \{f_{i,j}\}_{i,j\in I, i\leq j})$, where
  for all $i\leq j$, $f_{i,j} : H_p(K_i) \mapsto H_p(K_j)$ is induced
  by the inclusion map $K_i \mapsto K_j$.
\end{defn}

The functions $f_{i,j}$ allow us to link generators in each successive
homology space in the filtration. Since each generator correspond to a
topological feature (connected component, hole, void, etc, depending
on the dimension $p$), we can determine whether it survives in the
next step of the filtration. We can now determine when each feature is
born and when it dies (if it dies at all). This representation will be
dependent on the choice of basis for each homology space
$H_p(K_i)$. However, by the Fundamental Theorem of Persistent
Homology, we can choose base vectors in each homology space such that
the collection of half-open intervals is well-defined and unique. This
construction is called a \emph{barcode}.
%% TODO references for the Fundamental Theorem

\section{Topological summaries: barcodes and persistence diagrams}%
\label{sec:topol-summ}

In order to interpret the results of the persistent homology
computation, we need to compare the output for a particular data set
to a suitable null model. For this, we need some kind of a similarity
measure between barcodes and a way to evaluate the statistical
significance of the results.

One possible approach for this is to define a space in which we can
project barcodes and study their geometric
properties. \emph{Persistence diagrams} are an example of such a
space.

\begin{defn}[Persistence diagrams]
  A \emph{persistence diagram} is the union of a finite multiset of
  points in $\overline{\mathbb{R}}^2$ zith the diagonal
  $\Delta = \{(x,x) \;|\; x\in\mathbb{R}^2\}$, where every point of
  $\Delta$ has infinite multiplicity.
\end{defn}

The diagonal $\Delta$ is added to facilitate comparisons between
diagrams, as points near the diagonal correspond to short-lived
topological feature, thus likely to be caused by small perturbations
in the data.

We can now define several distances on the space of persistence
diagrams.

\begin{defn}[Wasserstein distance]
  The \emph{$p$-th Wasserstein distance} between two diagrams $X$ and
  $Y$ is
  \[ W_p[d](X, Y) = \inf_{\phi:X\mapsto Y} \left[\sum_{x\in X} {d\left(x, \phi(x)\right)}^p\right] \]
  for $p\in [1,\infty)$, and
  \[ W_\infty[d](X, Y) = \inf_{\phi:X\mapsto Y} \sup_{x\in X} d\left(x,
      \phi(x)\right) \] for $p = \infty$, where $d$ is a distance on
  $\mathbb{R}^2$ and $\phi$ ranges over all bijections from $X$ to
  $Y$.
\end{defn}

\begin{defn}[Bottleneck distance]
  The \emph{bottleneck distance} is defined as the infinite
  Wasserstein distance with $d$ the uniform norm:
  $d_B = W_\infty[L_\infty]$.
\end{defn}

Since the bottleneck distance is by far the most commonly used, we
will focus on it in the following. It is symmetric, non-negative, and
satisfies the triangle inequality. However, it is not a true distance,
as it is fairly straightforward to come up with two distinct diagrams
at bottleneck distance zero, even on multisets not touching the
diagonal $\Delta$.

\section{Stability}%
\label{sec:stability}

One of the most important aspects of Topological Data Analysis is that
it is \emph{stable} with respect to small perturbations in the
data. In fact, the persistence diagram operator is Lipschitz with
respect to the bottleneck distance. First, we define a distance
between subsets of a metric space.

\begin{defn}[Hausdorff distance]
  Let $X$ and $Y$ be subsets of a metric space $(E, d)$. The
  \emph{Hausdorff distance} is defined by
  \[ d_H(X,Y) = \max \left[ \sup_{x\in X} \inf_{y\in Y} d(x,y),
      \sup_{y\in Y} \inf_{x\in X} d(x,y) \right]. \]
\end{defn}

We can now give the proper stability property.

\begin{prop}
  Let $X$ and $Y$ be subsets in a metric space. We have
  \[ d_B(\dgm(X),\dgm(Y)) \leq d_H(X,Y). \]
\end{prop}

\section{Algorithms and implementations}%
\label{sec:algor-impl}

%% TODO
\cite{morozov_dionysus:_2018,bauer_ripser:_2018,reininghaus_dipha_2018,maria_gudhi_2014}

\chapter{Topological Data Analysis on Networks}%
\label{cha:topol-data-analys}

\section{Persistent homology for networks}%
\label{sec:pers-homol-netw}

We now consider the problem of applying persistent homology to network
data. An undirected network is already a simplicial complex of
dimension 1. However, this will not be sufficient to capture enough
topological information: we need to introduce higher-dimensional
simplices. The first possible method is to project the network on a
metric space~\cite{otter_roadmap_2017}, thus transforming the network
data into a point cloud data. For this, we need to compute the
distance between each pair of nodes in the network (via shortest path
distance for instance). This also requires the network to be
connected.

Another usual method for weighted networks is called the \emph{weight
  rank clique filtration} (WRCF)~\cite{petri_topological_2013}, which
filters the network based on weights. The procedure works as follows:
\begin{enumerate}
\item Set the set of all nodes, without any edge, as filtration
  step~0.
\item Rank all edge weights in decreasing order $\{w_1,\ldots,w_n\}$.
\item At filtration step $t$, keep only the edges whose weights are
  less than $w_t$, thus creating an unweighted graph.
\item Define the maximal cliques of the resulting graph to be
  simplices.
\end{enumerate}

At each step of the filtration, we construct a simplicial complex
based on cliques: this is called a \emph{clique complex}. It is
necessarily valid since a subset of a clique is necessarily a clique
itself, and the same is true for the intersection of two cliques.

This leads to a first possibility for applying persistent homology to
temporal networks. It is possible to segment the lifetime of the
network into sliding windows, creating a static graph on each window
by retaining only the edges available during the time interval. We can
then apply WRCF on each static graph in the sequence, obtaining a
filtered complex for each window, to which we can then apply
persistent homology.

This method can quickly become very computationally expensive, as
finding all maximal cliques (using the Bron-Kerbosch algorithm for
example) is a complicated problem in itself. In practice, we often
restrict the search to cliques of dimension lower than a certain bound
$d_M$. With this restriction, the new simplicial complex is
homologically equivalent to the original: they have the same homology
groups up to $H_{d_M-1}$.

This method is sensitive to the choice of sliding windows on the time
scale. The width and the overlap of the windows can completely change
the networks created and their topological features. Too small a
window, and the network becomes too small to have any significant
topological properties, too large, and we lose important information
in the evolution of the network over time.

\section{Zigzag persistence}%
\label{sec:zigzag-persistence}

The standard algorithm to compute persistent homology
(\autoref{sec:persistent-homology}) only works for filtrations which
are nested sequences of simplicial complexes:
\[ \cdots \subseteq K_{i-1} \subseteq K_i \subseteq K_{i+1} \subseteq
  \cdots \]

When studying temporal networks, we have two possibilities:
\begin{itemize}
\item Create an independent filtration (e.g.\ WRCF) from each time
  step. The issue is that the topological features will be completely
  disconnected from the time dimension.
\item Create a filtration along the time dimension. The issue in this
  case is that the sequence is no longer nested (except for additive
  temporal networks, ie when edges are never deleted).
\end{itemize}

The solution to consider the time dimension is provided by
\emph{zigzag persistence}~\cite{carlsson_zigzag_2009}, which allows to
compute persistence on alternating nested sequences:
\[ \cdots \supseteq K_{i-1} \subseteq K_i \supseteq K_{i+1} \subseteq
  \cdots \]

This sequence can in turn be computed from a temporal network by
computing the union of each pair of consecutive time steps,
constructing a alternating sequence.

Zigzag persistence is a special case of the more general concept of
\emph{multi-parameter
  persistence}~\cite{carlsson_theory_2009,dey_computing_2014}, where
filtrations can span across multiple parameters.

%% Note about libraries implementing zigzag persistence: Dionysus

\chapter{Persistent Homology for Machine Learning applications}%
\label{cha:pers-homol-mach}

The output of persistent homology is not directly usable by most
statistical methods. Barcodes and persistence diagrams, being a
multiset of points in $\overline{\mathbb{R}}^2$, are not elements of a
metric space in which we could perform statistical computations.

The distances between persistence diagrams defined
in~\autoref{sec:topol-summ} allow us to compare different
outputs. From a statistical perspective, it is possible to use a
generative model of simplicial complexes, and use a distance between
persistence diagrams to measure the similarity of our observations
with this null model~\cite{adler_persistent_2010}. This would
effectively define a metric space of persistence diagrams. It is even
possible to define some statistical summaries (means, medians,
confidence intervals) on these
spaces~\cite{turner_frechet_2014,munch_probabilistic_2015}.

The issue with this approach is that metric spaces do not offer enough
algebraic structure to be amenable to most machine learning
techniques. One of the most recent development in the study of
topological summaries has been to find mappings between the space of
persistence diagrams and Banach spaces.

\section{Vectorization methods}%
\label{sec:vect-meth}

\subsection{Persistence landscapes}

\subsection{Persistence images}

\subsection{Tropical and arctic semirings}

\section{Kernel-based methods}%
\label{sec:kernel-based-methods}

\subsection{Persistent scale-space kernel}

\subsection{Persistence weighted gaussian kernel}

\subsection{Sliced Wasserstein kernel}

\section{Comparison}%
\label{sec:comparison}




\backmatter%

% \nocite{*}
\printbibliography%

\end{document}



%%% Local Variables:
%%% mode: latex
%%% TeX-master: t
%%% End:
